\chapter{R and RStudio Fundamentals}
\label{chap:r-basics}

\chapterguide
  {No prior R programming is required.}
  {set up R/RStudio; use scripts, variables, assignment, object classes, arithmetic, and keyboard input.}

\section{R and RStudio setup}

Install \term{R} first, then \term{RStudio Desktop}. Verify the setup in the Console:

\begin{lstlisting}[style=dsr]
1 + 2
\end{lstlisting}

The result \texttt{3} confirms that RStudio can use the R runtime.

\begin{keyidea}
R executes the code; RStudio provides the development environment.
\end{keyidea}

\subsection{Console versus Script}

Use the Console for immediate evaluation and scripts for saved, reproducible work.

\begin{utpfig}[RStudio panes]{The four panes. Everything in this book happens in
one of them, and the split between the two on the left is the difference between
a result you can reproduce and one you cannot.}{fig:rstudio}
\begin{tikzpicture}[x=1cm,y=1cm]
\node[fcbox,minimum width=52mm,minimum height=20mm,align=left,
      fill=UTPPaleBlue,draw=UTPBlue] (src) at (0,1.15)
  {\textbf{Source} \ \tiny (top left)\\[2pt]
   \tiny Your \texttt{.R} script. Edited, saved, rerun.\\
   \tiny Ctrl+Enter sends the current line down.};
\node[fcbox,minimum width=52mm,minimum height=20mm,align=left,
      fill=UTPPaleGray,draw=UTPMuted] (con) at (0,-1.15)
  {\textbf{Console} \ \tiny (bottom left)\\[2pt]
   \tiny Runs immediately, keeps no record.\\
   \tiny Good for a quick test, bad for an assignment.};
\node[fcbox,minimum width=48mm,minimum height=20mm,align=left,
      fill=UTPPaleTeal,draw=UTPTeal] (env) at (5.6,1.15)
  {\textbf{Environment} \ \tiny (top right)\\[2pt]
   \tiny Every object currently in memory, with its\\
   \tiny class and value. The visual form of \texttt{ls()}.};
\node[fcbox,minimum width=48mm,minimum height=20mm,align=left,
      fill=UTPPaleOrange,draw=UTPOrange] (plt) at (5.6,-1.15)
  {\textbf{Files / Plots / Help} \ \tiny (bottom right)\\[2pt]
   \tiny Where a chart appears, and where \texttt{?paste}\\
   \tiny opens. Also shows the working directory.};

\draw[-{Stealth[length=2mm]},semithick,draw=UTPBlue]
  (src.south) -- node[fclabel,right]{Ctrl+Enter} (con.north);
\end{tikzpicture}
\end{utpfig}

\begin{pitfall}
Work typed straight into the Console is gone when the session ends. If a
result cannot be reproduced by running a saved script from top to bottom in a
clean session, it is not reproducible --- and the Environment pane still
showing the right answer is not evidence that it is.
\end{pitfall}

\begin{dstable}
\begin{tabularx}{\linewidth}{@{}p{3.0cm}X X@{}}
\toprule
\rowcolor{DSPaleBlue!92}
Area & What the lab uses it for & Example \\
\midrule
R Console & Immediate commands and quick tests. & \texttt{2 + 3} \\
R Script & Saved statements that can be edited and rerun. & \texttt{test.R} \\
\bottomrule
\end{tabularx}
\end{dstable}

\begin{workedexample}
The first saved script in the handout is:
\begin{lstlisting}[style=dsr]
# My first program in R Programming
myString <- "Hello, World!"
print(myString)
\end{lstlisting}
\texttt{\#} begins a comment, \texttt{<-} assigns a value, and \texttt{print()} displays it.
\end{workedexample}

\section{Packages}

Install a package once, then load it in each session that uses it:

\begin{lstlisting}[style=dsr]
install.packages("package_name")
library(package_name)
\end{lstlisting}

\texttt{install.packages()} installs; \texttt{library()} loads.

\section{Variables and object names}

Create variables and inspect the workspace:

\begin{lstlisting}[style=dsr]
var.1 = 5
var_1 = 7
x = 1
print(ls())
print(ls(pattern = "var"))
\end{lstlisting}

\texttt{ls()} lists workspace objects; its \texttt{pattern} argument filters their names. R object names may contain periods or underscores.

\section{Assignment operators}

R supports assignment in either direction:

\begin{dstable}
\begin{tabularx}{\linewidth}{@{}p{3.3cm}X@{}}
\toprule
\rowcolor{DSPaleBlue!92}
Operator & Handout description \\
\midrule
\texttt{<-}, \texttt{<<-}, \texttt{=} & Leftwards assignment \\
\texttt{->}, \texttt{->>} & Rightwards assignment \\
\bottomrule
\end{tabularx}
\end{dstable}

The examples combine assignment with several value types:

\begin{lstlisting}[style=dsr]
Gender <- "Female"
height <<- 152
Weight = 81
3 -> f
23.5 ->> x
b <- seq(from = 0, to = 5)
c <- seq(from = 0, to = 10, by = 2)
v = 2L
w <<- 2 + 5i
a <- charToRaw("Hello")
\end{lstlisting}

\texttt{seq()} generates sequences; \texttt{2L} is an integer; \texttt{2 + 5i} is complex; and \texttt{charToRaw()} returns raw byte values.

\begin{pitfall}
\texttt{<<-} and \texttt{->>} assign outside the current environment. Prefer \texttt{<-} for ordinary work because nonlocal assignment can make code difficult to trace.
\end{pitfall}

\subsection{Inspecting an object's class}

The handout checks object classes with \texttt{class()}:

\begin{lstlisting}[style=dsr]
print(class(Gender))
print(class(height))
print(class(f))
print(class(x))
print(class(b))
print(class(v))
print(class(w))
\end{lstlisting}

\begin{keyidea}
Use \texttt{class()} to inspect an object's class instead of inferring it from the name.
\end{keyidea}

\section{Arithmetic operations}

R provides the following arithmetic operators:

\begin{dstable}
\begin{tabularx}{\linewidth}{@{}p{2.2cm}X@{}}
\toprule
\rowcolor{DSPaleBlue!92}
Operator & Operation \\
\midrule
\texttt{+} & Addition \\
\texttt{-} & Subtraction \\
\texttt{*} & Multiplication \\
\texttt{**} or \texttt{\^} & Exponentiation \\
\texttt{/} & Division \\
\texttt{\%\%} & Modulus / remainder \\
\texttt{\%/\%} & Integer division \\
\bottomrule
\end{tabularx}
\end{dstable}

The examples combine these operators in a BMI calculation:

\begin{lstlisting}[style=dsr]
print(f + 3)
print(height - x)
print(Weight * 2)
print(b ** 2)
print(c ^ 5)

m = height / 100
print(Weight / (m ** 2))
BMI = Weight / (m ** 2)

print(b %% 2)
print(c %/% 2)
\end{lstlisting}

The BMI calculation reuses the height in meters and stores the result.

\section{Reading user input}

\texttt{readline()} returns text, so numeric input requires conversion:

\begin{lstlisting}[style=dsr]
name <- readline(prompt = "Enter name: ")
age <- readline(prompt = "Enter age: ")

# convert character into numeric
age <- as.numeric(age)

print(paste("Hi,", name,
            "this year you are", age, "years old."))
\end{lstlisting}

This is a basic input--process--output pattern.

\section{Built-in help and demonstrations}

Use built-in help and demonstrations to explore functions:

\begin{lstlisting}[style=dsr]
?paste
demo(graphics)
demo(image)
\end{lstlisting}


\begin{quickcheck}
\begin{enumerate}
  \item Why does the handout install R before RStudio?
  \item What is the practical difference between testing one expression in the Console and saving several statements in an R Script?
  \item What does \texttt{ls(pattern = "var")} do differently from \texttt{ls()}?
  \item Which operators in the handout return a remainder and an integer quotient?
  \item Why does the user-input example call \texttt{as.numeric(age)}?
\end{enumerate}
\end{quickcheck}

\begin{chapterrecap}
\begin{itemize}
  \item Use the Console for quick tests and scripts for reproducible programs.
  \item Prefer \texttt{<-} for routine assignment; inspect objects with \texttt{class()} and \texttt{ls()}.
  \item Arithmetic and input conversion support simple input--process--output programs.
\end{itemize}
\end{chapterrecap}
